% Fuentes: EXP-013, EXP-019..EXP-021; DEC-009.
\section{Conclusiones y siguientes pasos}
\begin{itemize}
  \item \textbf{Lo que funciona:} BatchNorm (con una \emph{learning rate} mayor) y la augmentation geométrica, que
  codifica la invarianza a la orientación de las imágenes aéreas; juntas suman +\res{gain-total}~pp de \emph{mean
  recall}.
  \item \textbf{Lo que no:} dropout y L2 bajan train y validación a la vez, y la variación fotométrica empeora.
  \item \textbf{Límite:} la brecha se reduce pero no desaparece; una red densa sobre píxeles no comparte pesos entre
  posiciones ni tiene invarianza a traslaciones~\cite{goodfellow}.
  \item \textbf{Siguiente paso:} redes convolucionales propias y arquitecturas conocidas entrenadas desde cero.
\end{itemize}

Todas las cifras proceden de \texttt{experiments/experiments.csv}; el razonamiento de cada decisión está en la
bitácora del proyecto.

\begin{thebibliography}{9}\small
\bibitem{xview} D. Lam et al. \emph{xView: Objects in Context in Overhead Imagery}. arXiv:1802.07856, 2018.
\bibitem{srivastava} N. Srivastava et al. \emph{Dropout: A Simple Way to Prevent Neural Networks from Overfitting}. JMLR 15, 2014.
\bibitem{ioffe} S. Ioffe, C. Szegedy. \emph{Batch Normalization: Accelerating Deep Network Training by Reducing Internal Covariate Shift}. ICML, 2015.
\bibitem{goodfellow} I. Goodfellow, Y. Bengio, A. Courville. \emph{Deep Learning}. MIT Press, 2016.
\bibitem{bouthillier} X. Bouthillier et al. \emph{Accounting for Variance in Machine Learning Benchmarks}. MLSys, 2021.
\bibitem{cawley} G. Cawley, N. Talbot. \emph{On Over-fitting in Model Selection and Subsequent Selection Bias in Performance Evaluation}. JMLR 11, 2010.
\end{thebibliography}
